% ===============================
% Template de Relatório de PoC (INE5448)
% IA & Segurança — LabREC/UFSC
% Versão 2.1
% ===============================
% INSTRUÇÕES GERAIS:
%   1. Substitua todos os campos marcados com [PREENCHER] pelas suas informações.
%   2. Mantenha a estrutura de seções — adicione subseções conforme necessário.
%   3. As seções marcadas com (*) são obrigatórias.
%   4. Remova os comentários de instrução antes de submeter.
% ===============================

\documentclass[12pt,a4paper]{article}

% ---- Codificação, idioma e fontes ----
\usepackage[utf8]{inputenc}
\DeclareUnicodeCharacter{2192}{\ensuremath{\rightarrow}}
\DeclareUnicodeCharacter{2190}{\ensuremath{\leftarrow}}
\DeclareUnicodeCharacter{2194}{\ensuremath{\leftrightarrow}}
\DeclareUnicodeCharacter{00D7}{\ensuremath{\times}}
\DeclareUnicodeCharacter{2265}{\ensuremath{\geq}}
\DeclareUnicodeCharacter{2264}{\ensuremath{\leq}}
\DeclareUnicodeCharacter{00B1}{\ensuremath{\pm}}
\usepackage[T1]{fontenc}
\usepackage[brazil]{babel}
\usepackage[a4paper,margin=2.5cm]{geometry}
\usepackage{times}

% ---- Matemática, figuras e links ----
\usepackage{amsmath,amssymb}
\usepackage{graphicx}
\usepackage{float}
\usepackage{hyperref}
\usepackage{enumitem}
\setlist[itemize]{leftmargin=1.5em}
\setlist[enumerate]{leftmargin=1.5em}
\usepackage{xcolor}
\usepackage{booktabs}
\usepackage{csquotes}
\usepackage{listings}
\usepackage{subcaption}
\usepackage{multirow}

% ---- Configuração de hyperlinks ----
\hypersetup{
    colorlinks=true,
    linkcolor=blue,
    filecolor=magenta,      
    urlcolor=cyan,
    citecolor=blue,
    pdftitle={Relatório de PoC - INE5448},
    pdfpagemode=FullScreen,
}

% ---- Configuração de listagens de código ----
\lstset{
    basicstyle=\ttfamily\small,
    breaklines=true,
    frame=single,
    numbers=left,
    numberstyle=\tiny,
    tabsize=2,
    commentstyle=\color{gray},
    keywordstyle=\color{blue},
    stringstyle=\color{red},
    literate=
        {á}{{\'a}}1 {é}{{\'e}}1 {í}{{\'i}}1 {ó}{{\'o}}1 {ú}{{\'u}}1
        {Á}{{\'A}}1 {É}{{\'E}}1 {Í}{{\'I}}1 {Ó}{{\'O}}1 {Ú}{{\'U}}1
        {à}{{\`a}}1 {è}{{\`e}}1 {ì}{{\`i}}1 {ò}{{\`o}}1 {ù}{{\`u}}1
        {À}{{\`A}}1 {È}{{\'E}}1 {Ì}{{\`I}}1 {Ò}{{\`O}}1 {Ù}{{\`U}}1
        {ä}{{\"a}}1 {ë}{{\"e}}1 {ï}{{\"i}}1 {ö}{{\"o}}1 {ü}{{\"u}}1
        {Ä}{{\"A}}1 {Ë}{{\"E}}1 {Ï}{{\"I}}1 {Ö}{{\"O}}1 {Ü}{{\"U}}1
        {â}{{\^a}}1 {ê}{{\^e}}1 {î}{{\^i}}1 {ô}{{\^o}}1 {û}{{\^u}}1
        {Â}{{\^A}}1 {Ê}{{\^E}}1 {Î}{{\^I}}1 {Ô}{{\^O}}1 {Û}{{\^U}}1
        {ã}{{\~a}}1 {õ}{{\~o}}1 {Ã}{{\~A}}1 {Õ}{{\~O}}1
        {ç}{{\c{c}}}1 {Ç}{{\c{C}}}1
        {ñ}{{\~n}}1 {Ñ}{{\~N}}1
}

\lstdefinelanguage{yaml}{
  keywords={true,false,null,y,n},
  keywordstyle=\color{blue}\bfseries,
  basicstyle=\ttfamily\small,
  sensitive=false,
  comment=[l]{\#},
  morecomment=[s]{/*}{*/},
  commentstyle=\color{gray}\ttfamily,
  stringstyle=\color{red}\ttfamily,
  moredelim=[l][\color{orange}]{\&},
  moredelim=[l][\color{magenta}]{*},
  moredelim=**[il][\color{green}]{:},
  morestring=[b]',
  morestring=[b]",
  literate=
    {á}{{\'a}}1 {é}{{\'e}}1 {í}{{\'i}}1 {ó}{{\'o}}1 {ú}{{\'u}}1
    {â}{{\^a}}1 {ê}{{\^e}}1 {î}{{\^i}}1 {ô}{{\^o}}1 {û}{{\^u}}1
    {ã}{{\~a}}1 {õ}{{\~o}}1 {ç}{{\c{c}}}1 {Ç}{{\c{C}}}1
}

\usepackage[table]{xcolor}

% ---- Bibliografia numérica (biber) ----
\usepackage[backend=biber,style=numeric,sorting=nyt,maxbibnames=99]{biblatex}
\addbibresource{ref.bib} 

% ---- Informações do documento ----
\title{\textbf{OCR-06: OCR Adaptativo com Aprendizagem Assistida}}
\author{Gabriel Salmoria \and Tutor: Lucas Mayr}

\date{\small \today}

% ==== Início do Documento ====
\begin{document}

\maketitle

% ---------------------------------------------------------------
% RESUMO (*)
% ---------------------------------------------------------------
\begin{abstract}
Sistemas de OCR baseados em Modelos de Linguagem Visual (VLMs) ainda enfrentam dificuldades na leitura de nomes próprios raros ou antigos. Frequentemente, o modelo tenta adivinhar uma palavra mais comum que ele encontrou durante o treinamento. Isso representa um problema sério em acervos e cartórios, onde os operadores precisam corrigir o texto manualmente. Além disso, o modelo genérico não retém essa correção para aprendizado futuro. O objetivo principal deste trabalho é desenvolver uma Prova de Conceito (PoC) focada na realização de fine-tuning em um VLM de código aberto utilizando LoRA. Para isso, o modelo é treinado de forma direcionada com recortes de linhas manuscritas do dataset BRESSAY e dados sintéticos cursivos de nomes atípicos (gerados a partir de bases do IBGE e TSE). Como complemento secundário a essa especialização do modelo, a PoC também explora a criação de um módulo de memória (via RAG e FAISS) para armazenar correções do usuário e mitigar falhas residuais. Este relatório descreve a definição inicial do problema, objetivos, escopo e o planejamento das próximas etapas, destacando as implicações de privacidade envolvidas no treinamento.
\end{abstract}

\textbf{Palavras-chave:} OCR, Modelos de Linguagem Visual, Fine-Tuning, LoRA, Privacidade, LGPD.

\tableofcontents

% ==============================================================
% SEÇÃO 1 — INTRODUÇÃO (*)
% ==============================================================
\section{Introdução}
\label{sec:introducao}

A digitalização de acervos históricos e registros públicos esbarra frequentemente nas limitações dos sistemas de OCR (Reconhecimento Ótico de Caracteres). Os modelos mais modernos conseguem processar bem o texto geral em português, mas apresentam alto índice de erro ao tratar de nomes próprios que fogem da distribuição padrão (a cauda longa). Como o modelo teve pouco ou nenhum contato com esses nomes durante o treinamento, ele tende a forçar uma "correção" para uma palavra mais comum do seu vocabulário \cite{cer_wer_ocr}. O resultado prático é a troca indevida de prenomes e sobrenomes, o que invalida a utilidade da transcrição no contexto de um Registro Civil e exige um longo esforço de revisão manual.

% --------------------------------------------------------------
\subsection{Motivação}

A principal motivação é que hoje o trabalho humano de correção é desperdiçado. O cartório paga alguém para consertar o texto gerado pelo OCR, mas o modelo não aprende nada com isso e continua cometendo o mesmo erro em documentos semelhantes. Historicamente, treinar um modelo gigante do zero exigiria infraestrutura proibitiva para instituições menores. No entanto, o uso de técnicas eficientes (como o LoRA) permite especializar o modelo diretamente no hardware disponível.

A proposta central deste trabalho é focar na melhoria direta do modelo por meio desse fine-tuning específico. Ao realizar um treinamento supervisionado (linha a linha) utilizando o dataset BRESSAY e dados cursivos artificiais baseados no IBGE e TSE, o objetivo é especializar a IA para a leitura correta do contexto manuscrito e de identidades da cauda longa. Apenas de forma complementar e secundária, o trabalho propõe acoplar uma "memória" baseada em busca semântica (RAG), servindo como auxílio extra na etapa de inferência para os casos de exceção que o modelo ajustado ainda não consiga acertar.

% --------------------------------------------------------------
\subsection{Objetivos}

O objetivo deste estudo é reduzir a taxa de erro de caracteres e palavras (CER e WER) na leitura de nomes atípicos em documentos manuscritos, buscando preservar a qualidade do OCR sobre o texto geral. Para alcançar essa meta, foram definidos os seguintes objetivos específicos:

\begin{itemize}
    \item Desenvolver e executar o fine-tuning de um VLM (via LoRA), utilizando anotações em nível de linha do dataset BRESSAY e um corpus sintético de nomes (baseado no IBGE e TSE), visando aumentar substancialmente a precisão do modelo em si.
    \item Implementar, de forma acessória, um banco vetorial (memória RAG) que utilize o histórico de correções do usuário como suporte secundário à predição.
    \item Analisar os ganhos de acurácia (CER e WER) proporcionados pelo fine-tuning e avaliar qual o benefício prático adicional obtido ao combinar essa estratégia principal com a memória episódica.
\end{itemize}

% --------------------------------------------------------------
\subsection{Descrição do Problema}

O problema central abordado neste projeto reside na incapacidade dos sistemas genéricos de OCR em reconhecer corretamente identidades atípicas (cauda longa) em documentos históricos e registros civis manuscritos. Como os modelos de linguagem visual (VLMs) tendem a ser enviesados pela distribuição do vocabulário comum, eles frequentemente distorcem prenomes e sobrenomes raros em tentativas de forçar uma coerência semântica. 

Como consequência, os operadores de acervos e cartórios são obrigados a realizar a correção manual das transcrições. Ao propormos uma solução baseada em fine-tuning com documentos reais para acabar com esse ciclo, esbarramos em um segundo problema de ordem ética e legal: o uso de registros civis autênticos para treinar modelos de inteligência artificial envolve o processamento de dados pessoais sensíveis (nomes reais e caligrafias autênticas), o que impõe desafios críticos de privacidade e conformidade estrita com a LGPD.

% ==============================================================
% SEÇÃO 2 — FUNDAMENTAÇÃO TEÓRICA (*)
% ==============================================================
\section{Fundamentação Teórica e Estado da Arte}
\label{sec:fundamentacao}

Atualmente, o processamento de OCR baseia-se fortemente em Modelos de Linguagem Visual (VLMs). Diferente dos métodos que operam estritamente no nível de caractere, essas arquiteturas processam a imagem em sua totalidade, extraindo a informação de maneira mais contextualizada.

% --------------------------------------------------------------
\subsection{Técnicas Relevantes para o Problema}

\subsubsection{Abordagens Clássicas}
Sistemas legados frequentemente utilizam o Tesseract ou arquiteturas compostas por redes LSTM. A principal limitação dessas abordagens mais tradicionais é a falta de compreensão de contexto de longo alcance. Consequentemente, elas tendem a falhar ao lidar com textos manuscritos contínuos, onde a sobreposição dos traços aumenta a ambiguidade.

\subsubsection{Abordagens Baseadas em ML/DL}
O estado da arte recente tem se concentrado no uso de Transformers, exemplificados pelo TrOCR \cite{cer_wer_ocr} e por modelos generativos robustos como a família Qwen-VL \cite{qwen2vl}. Para contornar os altos custos computacionais de treinar esses grandes modelos a partir do zero, o uso de técnicas de Parameter-Efficient Fine-Tuning (PEFT), como o LoRA \cite{lora2022}, permite adaptar os pesos utilizando hardwares comerciais mais acessíveis. Adicionalmente, pacotes como o Augraphy \cite{augraphy2022} têm sido empregados para gerar dados sintéticos que mimetizam manchas, ruídos e o envelhecimento natural do papel encontrado nos registros reais.

% --------------------------------------------------------------
\subsection{Explicabilidade em Modelos de IA}
[PREENCHER NA ENTREGA 2: Abordagem das técnicas de interpretabilidade e como os mapas de saliência das predições interagem com a interface de correção do usuário.]

% --------------------------------------------------------------
\subsection{Segurança e Privacidade em Sistemas de IA}

A adaptação de um VLM através de fine-tuning utilizando documentos que contêm dados reais expõe o sistema a riscos específicos de segurança e privacidade inerentes ao treinamento de modelos de linguagem.

\subsubsection{Extração de Dados e Memorização}
[PREENCHER NA ENTREGA 2: Detalhamento sobre como a IA pode acabar decorando nomes reais do conjunto de treinamento (overfitting) e o risco de vazamento dessas informações privadas durante a inferência.]

\subsubsection{Conformidade com a LGPD}
Por operar com fine-tuning sobre registros civis, o processo de treinamento precisa estar alinhado às exigências da LGPD. O uso do dataset BRESSAY e a geração de dados artificiais devem garantir que o modelo aprenda os traços da caligrafia e a estrutura de nomes atípicos, mas sem fixar vínculos com indivíduos reais. A estratégia central envolve a injeção forte de dados sintéticos (nomes cruzados aleatoriamente do IBGE e TSE) para diluir a presença de pessoas reais no treino, alinhando a ferramenta ao princípio de minimização e anonimização de dados. Adicionalmente, caso o módulo secundário de memória RAG seja ativado, ele armazenará estritamente "embeddings" irreversíveis, sem reter a imagem crua do documento.

% --------------------------------------------------------------
\subsection{Gaps e Oportunidades}

\begin{enumerate}
    \item \textbf{Falta de variabilidade nos datasets} — A base disponível de textos manuscritos carece de nomes próprios raros. Este projeto resolve essa limitação estruturando um pipeline de treinamento (fine-tuning) que incorpora um corpus sintético com dados do IBGE/TSE ao lado do dataset BRESSAY.
    \item \textbf{Baixa generalização para manuscritos complexos} — Modelos pré-treinados falham na leitura de caligrafias antigas. O foco no fine-tuning direcionado (linha a linha) ataca diretamente essa deficiência.
    \item \textbf{Custo de atualização (Gap Secundário)} — Como forma de suporte extra ao modelo especializado, a adição da memória baseada em RAG permite que correções pontuais feitas pelo usuário retroalimentem a inferência localmente, mitigando a necessidade técnica de re-treinamentos constantes para casos atípicos residuais.
\end{enumerate}

% ==============================================================
% SEÇÃO 3 — OBJETIVOS E MÉTRICAS DE SUCESSO (*)
% ==============================================================
\section{Objetivos e Métricas de Sucesso}
\label{sec:objetivos}

O sucesso deste trabalho depende de conseguir demonstrar que o modelo aprimorado melhora efetivamente o suporte humano, aprendendo a recuperar os termos difíceis de forma prática sem simplesmente memorizar o treinamento.

% --------------------------------------------------------------
\subsection{Objetivos do Estudo}

\begin{itemize}
    \item Estruturar um script gerador de imagens para criar um dataset robusto de texto cursivo sintético focado em nomes raros do IBGE e TSE, degradados artificialmente.
    \item Executar o fine-tuning intensivo de um VLM (via LoRA) utilizando a união de trechos alinhados do dataset BRESSAY e das imagens sintéticas geradas.
    \item Como suporte terciário, desenvolver e acoplar a arquitetura do banco vetorial FAISS como uma memória de correção e mensurar seu impacto prático nas predições do modelo ajustado.
\end{itemize}

% --------------------------------------------------------------
\subsection{Métricas de Sucesso}

\begin{enumerate}
    \item \textbf{CER (Character Error Rate) segmentado} 
    A métrica tradicional será avaliada de forma independente para nomes atípicos, nomes comuns e texto em geral, assegurando que o fine-tuning traga ganhos em domínios específicos sem prejudicar o desempenho geral.

    \item \textbf{Hit Rate da Memória} 
    Proporção de vezes em que o uso acessório do histórico do usuário permitiu ao modelo reverter uma predição isoladamente incorreta.

    \item \textbf{[PREENCHER: Métrica 3]} 
    [PREENCHER NA ENTREGA 2]
\end{enumerate}

% ==============================================================
% SEÇÃO 4 — ESCOPO DA SOLUÇÃO E REQUISITOS (*)
% ==============================================================
\section{Escopo da Solução e Requisitos}
\label{sec:escopo}

A solução proposta será entregue em formato de backend, estruturado para expor serviços de processamento de OCR refinados pelo novo treinamento.

% --------------------------------------------------------------
\subsection{Escopo da Solução}

O escopo do projeto inclui:
\begin{itemize}
    \item O pré-processamento de imagens do dataset BRESSAY e os scripts responsáveis por gerar/corromper o dataset sintético de nomes (baseado em IBGE/TSE).
    \item A configuração e execução primária do fine-tuning no VLM escolhido, assim como seu fluxo de inferência.
    \item A integração de um banco vetorial (como o FAISS) para atuar de forma auxiliar como memória do histórico de correções.
\end{itemize}

Estão fora do escopo inicial:
\begin{itemize}
    \item O desenvolvimento de interfaces gráficas completas para o usuário final.
    \item A integração e implantação da ferramenta em um sistema legado real de cartório.
    \item O processamento voltado para idiomas estrangeiros ou para textos inteiramente tipográficos.
\end{itemize}

% --------------------------------------------------------------
\subsection{Requisitos Funcionais}

\begin{enumerate}
    \item O sistema deve receber um recorte de imagem manuscrita e retornar a transcrição textual baseada no modelo especializado.
    \item O ambiente de fine-tuning deve permitir a parametrização do LoRA e a experimentação iterativa rápida.
    \item Deve existir um fluxo (API) para que as correções do usuário sejam encaminhadas ao sistema.
    \item Durante a inferência, a rotina de predição pode opcionalmente consultar o banco de memória antes de determinar a resposta final.
\end{enumerate}

% --------------------------------------------------------------
\subsection{Requisitos Não-Funcionais}

\begin{itemize}
    \item \textbf{Desempenho:} O treinamento do modelo (fine-tuning) precisa ser comportado por uma GPU comercial de até 24 GB de VRAM, eliminando a dependência de clusters de altíssimo custo.
    \item \textbf{Reprodutibilidade:} O particionamento dos dados, as seeds de pseudo-aleatoriedade e as configurações de hiperparâmetros devem estar rigorosamente documentadas no repositório.
    \item \textbf{Segurança e Privacidade:} O pipeline de fine-tuning deve garantir que a predominância dos dados atípicos seja sintética (IBGE/TSE) para mitigar a memorização de dados de cidadãos reais presentes nas amostras do BRESSAY.
\end{itemize}

% ==============================================================
% SEÇÃO 5 — AMEAÇAS, RISCOS E CONTROLES (*)
% ==============================================================
\section{Ameaças, Riscos e Controles}
\label{sec:threat-model}

O treinamento (fine-tuning) de modelos de inteligência artificial com dados baseados em registros públicos cria potenciais superfícies de ataque relacionadas à privacidade dos dados de treino.

% --------------------------------------------------------------
\subsection{Principais Ameaças Identificadas}

\begin{itemize}
    \item \textbf{Ataques de Extração de Dados (Data Extraction):} Cenário onde um atacante conseguiria usar prompts específicos para forçar o modelo a "cuspir" nomes reais e combinações de dados sensíveis que a rede memorizou durante o fine-tuning.
    \item \textbf{Ataques de Inversão de Modelo:} A tentativa de inferir se um determinado documento ou nome de cidadão específico fez parte do dataset de treinamento do modelo (Membership Inference).
    \item \textbf{Poisoning (Envenenamento de Dados):} Inserções acidentais ou maliciosas de anotações incorretas no dataset de treino (ou no módulo de memória secundário), fazendo o modelo reproduzir erros sistêmicos em produção.
    \item \textbf{[PREENCHER: ameaça 4]:} [PREENCHER NA ENTREGA 2]
\end{itemize}

% --------------------------------------------------------------
\subsection{Riscos Associados ao Projeto}

\begin{enumerate}
    \item \textbf{Restrição de Hardware:} O risco de o limite de VRAM de 24 GB impedir um fine-tuning de qualidade adequada, obrigando a aplicação de métodos de quantização agressiva que prejudiquem a performance final.
    \item \textbf{Viés de Representatividade:} O risco de os dados sintéticos gerados (via IBGE/TSE) não representarem uniformemente as origens demográficas, ocasionando viés preditivo em caligrafias ou regiões específicas.
    \item \textbf{[PREENCHER: risco 3]:} [PREENCHER NA ENTREGA 2]
\end{enumerate}

% --------------------------------------------------------------
\subsection{Controles e Medidas de Mitigação}

\begin{itemize}
    \item Injeção de volume majoritário de dados sintéticos aleatórios (IBGE/TSE) no fine-tuning para evitar que o modelo sofra overfitting nos poucos dados reais do BRESSAY, reduzindo substancialmente o risco de memorização e extração.
    \item Instituir uma separação rigorosa entre os conjuntos de dados de treinamento, validação e teste para evitar o vazamento de informações ("data leakage") e garantir a avaliação íntegra da capacidade de generalização.
    \item Para o módulo auxiliar (RAG), armazenar de forma persistente apenas os "embeddings" não-reversíveis extraídos pelo modelo, descartando qualquer imagem original do documento.
\end{itemize}

% ==============================================================
% SEÇÃO 6 — METODOLOGIA DE DESENVOLVIMENTO DA PoC (*)
% ==============================================================
\section{Metodologia de Desenvolvimento da PoC}
\label{sec:metodologia}

Os experimentos farão uso conjunto de recortes provenientes do corpus BRESSAY e de amostras manuscritas inteiramente sintéticas geradas pelo projeto. O conjunto sintético baseia-se em cruzamentos das listas de prenomes do Censo IBGE de 2010 com sobrenomes extraídos de registros públicos da Receita Federal e do Tribunal Superior Eleitoral. Por meio da combinação dessas bases, foi estruturado um pipeline capaz de montar nomes completos de alta atipicidade e convertê-los em imagens de caligrafias simuladas para fortalecer o fine-tuning.

% --------------------------------------------------------------
\subsection{Fases de Implementação}

\begin{enumerate}
    \item \textbf{Fase 1: Preparação dos Dados} — Filtragem das bases do IBGE/Receita, estruturação do gerador de caligrafia sintética e aplicação da biblioteca Augraphy para adicionar artefatos visuais ao papel.
    \item \textbf{Fase 2: Fine-Tuning do VLM} — Configuração de um modelo de linguagem visual (como o Qwen-VL) integrado à biblioteca Unsloth, visando a aceleração e otimização do treinamento (via LoRA) sobre a base consolidada de recortes manuscritos.
    \item \textbf{Fase 3: Construção da Memória} — Parametrização do banco FAISS como suporte de inferência para o gerenciamento secundário de vetores de correções.
    \item \textbf{Fase 4: Validação} — Apuração final e comparativa das métricas de CER e WER focando nas predições do modelo ajustado isoladamente, e posteriormente com o auxílio do RAG.
\end{enumerate}

% --------------------------------------------------------------
\subsection{Pipelines / Módulos Desenvolvidos}

\begin{itemize}
    \item \textbf{Módulo de Síntese:} Responsável por ler os dados estruturados de nomes, definir tipografias variadas e construir a massa de dados (imagens emparelhadas aos textos) que servirá de treinamento principal ao lado do BRESSAY.
    \item \textbf{Módulo de Inferência e RAG:} O código backend final, responsável por unir os pesos-base do modelo com o pacote LoRA recém-treinado e, opcionalmente, fazer a orquestração do acesso ao banco FAISS.
\end{itemize}

% --------------------------------------------------------------
\subsection{Ferramentas e Tecnologias Utilizadas}

\begin{itemize}
    \item \textbf{Python e PyTorch:} O ambiente fundamental de desenvolvimento para orquestração e treino dos modelos.
    \item \textbf{Unsloth e ecossistema HuggingFace:} A adoção da biblioteca Unsloth auxilia substancialmente na redução do consumo de VRAM, acelerando a fase de fine-tuning e viabilizando o projeto em hardware pessoal.
    \item \textbf{FAISS:} Ferramenta otimizada desenvolvida para viabilizar buscas semânticas rápidas, formando o mecanismo de memória auxiliar do sistema.
    \item \textbf{Git e DVC:} Utilizados para estruturar o versionamento seguro tanto da base de código (Github) quanto dos datasets primários (Data Version Control).
\end{itemize}

% --------------------------------------------------------------
\subsection{Critérios de Aceitação da PoC}

\begin{enumerate}
    \item Geração comprovada de partições consistentes entre os datasets de treino e validação, garantindo a avaliação justa de nomes inéditos na fase final.
    \item A conclusão rotineira dos scripts de fine-tuning sem apresentar esgotamento de memória (Out-Of-Memory) no ambiente laboratorial disponibilizado.
    \item Demonstração de queda métrica na proporção de erros (CER/WER) nos nomes da cauda longa apenas com o uso do modelo fine-tunado, comprovando sua eficácia intrínseca.
    \item Validação de que o modelo não sofreu overfitting prejudicial, de modo que dados sensíveis autênticos não possam ser extraídos através da inferência.
\end{enumerate}

\subsection{Planejamento Inicial}

O cronograma segue o balizamento das semanas letivas restantes. Nas primeiras duas semanas o foco concentra-se puramente na geração refinada dos datasets sintéticos e integração do Augraphy. Após esta primeira entrega, as tarefas voltam-se primariamente para a codificação da rotina de fine-tuning do modelo e verificação exaustiva de sua convergência. O planejamento direciona o meio de outubro para a avaliação do modelo isolado e a implementação do suporte do FAISS. Novembro está reservado para a avaliação metódica dos resultados completos, finalizando os relatórios experimentais.

% ==============================================================
% SEÇÃO 7 — AVALIAÇÃO EXPERIMENTAL (*)
% ==============================================================
\section{Avaliação Experimental}
\label{sec:avaliacao}

[PREENCHER NA ENTREGA 2: parágrafo introdutório apresentando os cenários experimentais]

% --------------------------------------------------------------
\subsection{Experimento 1 — [PREENCHER: nome/descrição]}

[PREENCHER NA ENTREGA 2: descrição do setup experimental]

\begin{itemize}
    \item{} [PREENCHER: o que foi medido/coletado 1]
    \item{} [PREENCHER: o que foi medido/coletado 2]
\end{itemize}

[PREENCHER NA ENTREGA 2: análise dos resultados do experimento 1]

% --------------------------------------------------------------
\subsection{Experimento 2 — [PREENCHER: nome/descrição]}

[PREENCHER NA ENTREGA 3/4: descrição do setup experimental]

\begin{itemize}
    \item{} [PREENCHER: o que foi medido/coletado 1]
    \item{} [PREENCHER: o que foi medido/coletado 2]
\end{itemize}

[PREENCHER NA ENTREGA 3/4: análise dos resultados do experimento 2]

% ==============================================================
% SEÇÃO 8 — DISCUSSÃO (*)
% ==============================================================
\section{Discussão}
\label{sec:discussao}

[PREENCHER NA ENTREGA 4: parágrafo introdutório da discussão]

% --------------------------------------------------------------
\subsection{Alcance dos Objetivos}
[PREENCHER NA ENTREGA 4]

% --------------------------------------------------------------
\subsection{Limitações Identificadas}

\subsubsection{Limitações Técnicas}
[PREENCHER NA ENTREGA 4]

\subsubsection{Limitações de Segurança}
[PREENCHER NA ENTREGA 4]

% --------------------------------------------------------------
\subsection{Comparação com o Estado da Arte}
[PREENCHER NA ENTREGA 4]

% --------------------------------------------------------------
\subsection{Trade-offs Segurança vs. Usabilidade}
[PREENCHER NA ENTREGA 4]

% --------------------------------------------------------------
\subsection{Implicações Práticas}

\subsubsection{Para Usuários e Stakeholders}
[PREENCHER NA ENTREGA 4]

\subsubsection{Para Políticas Públicas}
[PREENCHER NA ENTREGA 4]

% ==============================================================
% SEÇÃO 9 — CONCLUSÕES E PRÓXIMOS PASSOS (*)
% ==============================================================
\section{Conclusões e Próximos Passos}
\label{sec:conclusoes}

[PREENCHER NA ENTREGA 4: parágrafo de abertura das conclusões]

% --------------------------------------------------------------
\subsection{Síntese das Evidências}

\begin{itemize}
    \item{} [PREENCHER: achado principal 1]
    \item{} [PREENCHER: achado principal 2]
\end{itemize}

% --------------------------------------------------------------
\subsection{Contribuições}

\begin{enumerate}
    \item{} [PREENCHER: contribuição 1]
\end{enumerate}

% --------------------------------------------------------------
\subsection{Trabalhos Futuros}

\subsubsection{Curto Prazo (3–6 meses)}
\begin{itemize}
    \item{} [PREENCHER]
\end{itemize}

\subsubsection{Médio Prazo (6–12 meses)}
\begin{itemize}
    \item{} [PREENCHER]
\end{itemize}

\subsubsection{Longo Prazo (1–2 anos)}
\begin{itemize}
    \item{} [PREENCHER]
\end{itemize}

% --------------------------------------------------------------
\subsection{Recomendações}
[PREENCHER NA ENTREGA 4]

% ==============================================================
% SEÇÃO 10 — CHECKLIST DE CONFORMIDADE (*)
% ==============================================================
\section{Checklist de Conformidade (LGPD / Ética / Segurança)}
\label{sec:checklist}

[PREENCHER NA ENTREGA 2/3: parágrafo introdutório da seção]

% --------------------------------------------------------------
\subsection{Lei Geral de Proteção de Dados (LGPD)}

\begin{itemize}
    \item \textbf{Minimização de dados (Art. 6º, III)} \\
    Atendimento: [PREENCHER NA ENTREGA 2/3] \\
    Lacuna: [PREENCHER NA ENTREGA 2/3]

    \item \textbf{Finalidade específica (Art. 6º, I)} \\
    Atendimento: [PREENCHER NA ENTREGA 2/3] \\
    Lacuna: [PREENCHER NA ENTREGA 2/3]

    \item \textbf{Transparência (Art. 6º, VI)} \\
    Atendimento: [PREENCHER NA ENTREGA 2/3] \\
    Lacuna: [PREENCHER NA ENTREGA 2/3]

    \item \textbf{Segurança (Art. 46)} \\
    Atendimento: [PREENCHER NA ENTREGA 2/3] \\
    Lacuna: [PREENCHER NA ENTREGA 2/3]

    \item \textbf{Não discriminação (Art. 6º, IX)} \\
    Atendimento: [PREENCHER NA ENTREGA 2/3] \\
    Lacuna: [PREENCHER NA ENTREGA 2/3]
\end{itemize}

% --------------------------------------------------------------
\subsection{Segurança da Informação}

\begin{itemize}
    \item \textbf{Acesso e Controle} \\
    Atendimento: [PREENCHER NA ENTREGA 2/3] \\
    Lacuna: [PREENCHER NA ENTREGA 2/3]

    \item \textbf{Integridade e Reprodutibilidade} \\
    Atendimento: [PREENCHER NA ENTREGA 2/3] \\
    Lacuna: [PREENCHER NA ENTREGA 2/3]

    \item \textbf{Proteções Técnicas} \\
    Atendimento: [PREENCHER NA ENTREGA 2/3] \\
    Lacuna: [PREENCHER NA ENTREGA 2/3]
\end{itemize}

% --------------------------------------------------------------
\subsection{Ética e IA Responsável}

\begin{itemize}
    \item \textbf{Transparência Experimental} \\
    Atendimento: [PREENCHER NA ENTREGA 2/3] \\
    Lacuna: [PREENCHER NA ENTREGA 2/3]

    \item \textbf{Responsabilidade} \\
    Atendimento: [PREENCHER NA ENTREGA 2/3] \\
    Lacuna: [PREENCHER NA ENTREGA 2/3]

    \item \textbf{Impacto Social} \\
    Atendimento: [PREENCHER NA ENTREGA 2/3] \\
    Lacuna: [PREENCHER NA ENTREGA 2/3]
\end{itemize}

% --------------------------------------------------------------
\subsection{Ações Necessárias para Maturidade em Produção}

\begin{itemize}
    \item{} [PREENCHER NA ENTREGA 2/3: ação prioritária 1]
    \item{} [PREENCHER NA ENTREGA 2/3: ação prioritária 2]
\end{itemize}

% ==============================================================
% AGRADECIMENTOS (opcional)
% ==============================================================
% \section*{Agradecimentos}
% [PREENCHER — ou remova esta seção se não aplicável]

% ==============================================================
% BIBLIOGRAFIA (*)
% ==============================================================
\nocite{*}
\printbibliography

% ==============================================================
% APÊNDICES (opcionais)
% ==============================================================
\appendix

% --------------------------------------------------------------
\section{Detalhes de Implementação}
\label{app:implementacao}

[PREENCHER NA ENTREGA 2/3: descrição da arquitetura de implementação]

\subsection{Estrutura de Diretórios}

\begin{verbatim}
[PREENCHER NA ENTREGA 2/3: cole a saída de `tree` ou descreva manualmente]
\end{verbatim}

\subsection{Comandos para Reprodução}

\subsubsection*{1) Instalação de Dependências}
\begin{verbatim}
[PREENCHER NA ENTREGA 2/3: comandos de instalação]
\end{verbatim}

\subsubsection*{2) Execução do Experimento Principal}
\begin{verbatim}
[PREENCHER NA ENTREGA 2/3: comando de execução]
\end{verbatim}

% --------------------------------------------------------------
\section{Exemplos Adicionais de Resultados}
\label{app:exemplos}

[PREENCHER NA ENTREGA 3/4: descrição dos exemplos apresentados neste apêndice]

% --------------------------------------------------------------
\section{Código-fonte}
\label{app:codigo}

\url{https://github.com/gabriel-salmoria/ai-sec}

\end{document}
